% !TeX root = ../main-cn.tex
% UTF-8. Abstract synchronized with paper/论文草稿.md.

文本到人体动作生成能够根据自然语言产生多样动作，但文本难以指定具体运动实例；稀疏惯性测量单元（IMU）提供连续的局部运动线索，却不足以唯一确定全身动作。
本文研究 IMU-guided Text-to-Motion Generation，以文本表达动作意图，以少量 IMU 读数约束生成动作的运动细节。
我们提出面向预训练文本到动作模型的轻量条件适配框架 ITM，将稀疏 IMU 表示、时序控制编码和生成基座适配组织为统一流程。
框架冻结生成骨干与文本编码器，通过可训练的控制模块注入 IMU 信息，以支持向不同生成架构迁移，并以保持文本语义和动作自然性为目标实现实例级运动控制。
我们构建 HumanML3D/AMASS 文本、动作与虚拟 IMU 配对数据，并通过扩散与流匹配基座上的实验考察框架的适用性。
受控条件对照检验 IMU 是否使生成动作更贴近目标运动，文本编辑检验语言能否补充未观测部位的语义，生成质量评价则衡量新增控制对语义匹配和动作分布的影响。
实验揭示了局部运动约束与全身语义生成的互补潜力，以及控制精度、动作平滑性和语义保持之间的权衡。
本研究为将可穿戴传感器接入预训练动作生成器提供了可扩展的适配思路与评价依据。
% 作者注：正文采用面向多基座迁移的设计目标；“支持迁移”不等于任意基座无需适配，更不代表已在最新最强模型上验证。
% 当前已验证范围：MDM、MotionLab；HY-Motion 仅完成 Text-only smoke 和旋转导出，尚无融合训练与评价结果。
% 严谨结果备忘：MDM 多种子结果支持双腕实例级控制，单头收益不显著；不得将该结论推广为所有配置和基座均可靠。
% 严谨结果备忘：摆臂和转向有响应，速度与步幅响应不可靠；控制收益伴随 FID 代价与部分文本补全能力下降。
% 待补：HY-Motion 融合训练、paired/shuffled/zero/Text-only 对照，以及同协议生成质量和多种子控制评价。
% 待补：统一测试 ID、caption、物理时长和重复次数；不同基座的噪声空间不同，仅在同一基座内共享初始噪声。
% 待补：文本编辑同时比较 Text-only 与 Text+IMU 的语义响应和局部约束，检验互补性而非仅检验输出变化。
% 待补：报告各基座适配接口、训练参数和迁移成本；只有具备相应证据后，才将“设计目标”升级为“跨基座稳定有效”的实证结论。
